\section{The Road to October: Stalemate After 1967}

\subsection{From the Six-Day War to the War of Attrition}
The crushing Israeli victory in the Six-Day War of June 1967 redrew the map of the Middle East and left Israel occupying the Sinai Peninsula, the Golan Heights, and the West Bank. For Egypt the loss of Sinai was a national humiliation that demanded redress. Yet the diplomatic landscape froze almost immediately, as neither side could translate battlefield results into a durable settlement. Israeli confidence hardened into complacency, while Arab leaders searched for a way to reverse the verdict imposed upon them.

Between 1969 and 1970 the Egyptian-Israeli front along the Suez Canal erupted into the grinding War of Attrition. Gamal Abdel Nasser sought to wear down Israeli forces through artillery duels, commando raids, and air engagements, hoping that steady casualties would force political movement. The fighting drew in Soviet advisers and air-defense systems, raising the stakes of superpower involvement. A cease-fire in August 1970 ended the open conflict, but it resolved none of the underlying questions about Sinai's future or Egypt's lost territory.

Nasser's death in September 1970 brought Anwar Sadat to power, an outcome many observers initially dismissed as transitional. Sadat inherited a stalemate that was politically intolerable: Egypt could neither accept the post-1967 status quo nor break it through diplomacy alone. The frozen front lines symbolized a wider paralysis in which United Nations resolutions and great-power proposals produced no tangible withdrawal. Israeli leaders, secure behind the Bar-Lev Line, saw little reason to make concessions while their military superiority appeared overwhelming and unchallenged.

The diplomatic deadlock deepened through the early 1970s as repeated initiatives collapsed. Sadat declared 1971 a year of decision, then watched it pass without result, eroding his domestic credibility. Washington showed limited urgency, preoccupied with Vietnam and the broader contours of détente with Moscow. Israel, governed by Golda Meir and advised by Moshe Dayan, interpreted Arab inaction as evidence that deterrence held. This complacency, rooted in the memory of 1967, would prove a dangerous foundation for the assessments that followed.

By 1973 the structural conditions for renewed war were firmly in place. Egypt faced an unbearable choice between permanent loss and forceful action, while Israel underestimated the willingness of its adversaries to absorb costs for political ends. The War of Attrition had demonstrated that Arab armies could fight tenaciously, yet Israeli planners discounted the lesson. The stalemate after 1967 was thus not stability but a pressurized impasse, one that Sadat increasingly viewed as solvable only through a calculated and limited military shock. \cite{bregman2002} \cite{siniver2013} \cite{sadat1978}

\subsection{Sadat's Strategic Calculus and Limited War Aims}
Anwar Sadat approached the question of war not as a general seeking conquest but as a statesman seeking leverage. He understood that Egypt could not defeat Israel outright or reconquer the entire Sinai through force of arms. Instead, he framed military action as a means to shatter the diplomatic freeze and compel the superpowers to engage. A limited but credible offensive, he reasoned, would restore Egyptian bargaining power and national dignity, transforming the political conversation that had stagnated since the 1967 catastrophe.

Central to Sadat's calculus was the psychological dimension of the 1967 defeat. The humiliation had paralyzed Egyptian diplomacy, since no leader could negotiate from a position of evident weakness. Sadat believed that even a partial crossing of the Suez Canal, holding a strip of recaptured ground, would alter how Egyptians and the world perceived their standing. By restoring a sense of parity, he hoped to enter future talks as an equal rather than a supplicant pleading for the return of occupied land.

Sadat's planning therefore emphasized achievable objectives over sweeping ambitions. The aim was to seize and hold a defensible position on the east bank, inflict significant losses, and force international intervention before Israel could mount a decisive counterstroke. He coordinated closely with Syria, whose ambitions on the Golan Heights were more territorial, while keeping his own goals deliberately constrained. This divergence in war aims between the two allies would later shape the uneven and sometimes contradictory course of the fighting on the two fronts.

The strategic gamble also depended on managing the superpowers. Sadat expelled thousands of Soviet advisers in 1972, partly to assert independence and partly to signal flexibility toward Washington. He calculated that a sudden crisis would draw the United States and Soviet Union into mediation, since neither could tolerate an open-ended regional war threatening détente. By making the conflict dangerous enough to demand resolution, Sadat intended to convert military risk into diplomatic opportunity, a bet that few of his contemporaries fully appreciated.

In retrospect, Sadat's limited-war concept reveals a sophisticated fusion of force and diplomacy. He did not expect total victory and arguably did not require it; the political value lay in breaking the perception of permanent Egyptian impotence. Critics debate whether his intentions were as restrained as he later claimed in his memoirs, or whether ambitions expanded with early success. Regardless, the framing of war as a springboard to negotiation distinguishes 1973 from earlier rounds and helps explain its eventual diplomatic harvest. \cite{sadat1978} \cite{siniver2013} \cite{rabinovich2004}

\begin{figure}[htbp]
\centering
\includegraphics[width=0.88\linewidth,height=0.58\textheight,keepaspectratio]{figures/ch01_web.jpg}
\caption{Egyptian President Anwar Sadat, who pursued a limited war in 1973 to break the diplomatic stalemate and reopen negotiations over the Sinai.}
\label{fig:ch_egyptian_president_anwar_sadat_w}
\end{figure}

\bigskip
\begin{otherlanguage}{hebrew}
\subsection*{סיכום הפרק}

לאחר מלחמת ששת הימים נוצר קיפאון מדיני וצבאי בין ישראל למצרים. סאדאת תכנן מלחמה מוגבלת שמטרתה לשבור את הקיפאון ולכפות משא ומתן מחדש על סיני.

\end{otherlanguage}

\clearpage
